\documentclass[11pt]{article}

\usepackage[margin=1in]{geometry}
\usepackage[T1]{fontenc}
\usepackage[utf8]{inputenc}
\usepackage{lmodern}
\usepackage{microtype}

\usepackage{amsmath}
\usepackage{amssymb}
\usepackage{booktabs}
\usepackage{graphicx}
\usepackage{hyperref}
\usepackage{url}
\usepackage{xcolor}

\title{Convolutional Context Propagation for Long-Context Reasoning}
\author{\texttt{TODO}}
\date{}

\begin{document}
\maketitle

\begin{abstract}
This paper examines whether a fixed-topology, convolution-inspired pipeline can improve long-context reasoning while keeping inference cost predictable. We introduce Convolutional Context Propagation (CCP), which chunks long documents into overlapping segments, extracts local evidence nodes per chunk, and then performs sliding-window fusion over multiple layers to progressively abstract and consolidate information. A final verification step checks support and issues a corrected answer when the fused draft is weak. We compare CCP to a recursive baseline that iteratively expands a query into subproblems and composes intermediate answers. Under a controlled benchmark setup with a consistent evaluation protocol, CCP achieves higher judged answer quality while reducing estimated per-example cost and lowering the maximum number of model calls per example. These results suggest that hierarchical windowed aggregation can provide a practical accuracy--cost trade-off, with tighter cost predictability than recursive reasoning on long-context tasks.
\end{abstract}

\section{Introduction}

TODO: Motivate long-context failure modes (distraction, retrieval errors, cost blowups).
State the thesis: fixed-topology convolution-style aggregation can improve predictability.


\section{Method}

TODO: Define CCP formally.

\subsection{Chunking}
TODO: Define chunk length and overlap.

\subsection{Local Extraction}
TODO: Describe per-chunk node schema and evidence anchoring.

\subsection{Windowed Fusion (Convolution)}
TODO: Describe window size, stride, number of layers, and how nodes propagate upward.

\subsection{Verification}
TODO: Describe support check and correction behavior.


\section{Experiments}

TODO: Datasets, metrics (judge protocol), and cost accounting.

\subsection{Setup}
TODO: Model, decoding settings, and budgets.

\subsection{Results}
TODO: Tables/plots comparing CCP vs recursive baseline.


\section{Related Work}

TODO: Long-context methods, hierarchical summarization, recursive reasoning baselines.


\section{Limitations}

TODO: Judge bias, dataset coverage, verification brittleness, window hyperparameter sensitivity.


\section{Conclusion}

TODO: Summarize findings and practical implications.



\bibliographystyle{plain}
\bibliography{references}

\end{document}

